% Generated draft: requires mathematical and source review.
\documentclass[11pt]{article}
\usepackage[margin=25mm]{geometry}
\usepackage{fontspec}
\setmainfont{Latin Modern Roman}
\usepackage{amsmath,amssymb,mathtools,unicode-math}
\setmathfont{Latin Modern Math}
\usepackage{enumitem,needspace,xurl,hyperref}
\setlength{\parindent}{0pt}
\setlength{\parskip}{5pt}
\setlength{\emergencystretch}{3em}
\newcommand{\N}{\mathbb{N}}
\newcommand{\Z}{\mathbb{Z}}
\newcommand{\Q}{\mathbb{Q}}
\newcommand{\R}{\mathbb{R}}
\newcommand{\C}{\mathbb{C}}
\newcommand{\F}{\mathbb{F}}
\newcommand{\A}{\mathbb{A}}
\newcommand{\PP}{\mathbb{P}}
\newcommand{\E}{\mathbb{E}}
\newcommand{\eps}{\varepsilon}
\newcommand{\dd}{\,\mathrm{d}}
\begin{document}
\section*{UnsolvedMath ID: 30000273}
\textbf{Source problem code:} catalogue entry
\subsection{Definitions and mathematical statement}
\textbf{Characterization problem.} Characterize the integer sequences \((u_n)_{n\ge0}\) for which the ordinary generating function \(U(z)=\sum_{n\ge0}u_nz^n\) is algebraic. Here \(u_n\in\Z\), and under the natural characteristic-zero formal interpretation, algebraic means algebraic over \(\Q(z)\): there is a nonzero polynomial \(P(z,y)\in\Q[z,y]\), equivalently one in \(\Z[z,y]\) after clearing denominators, with\[P(z,U(z))=0\quad\text{in }\Q[[z]].\] The variable \(z\) is formal; no convergence-radius assumption, positivity, recurrence, or special formula for the coefficients is supplied. This is an ordinary, not exponential, generating function.\par The original statement asks to characterize the relation between integrality of all \(u_n\) and algebraicity of a generating function \(u\). The curated statement fixes integer coefficients first and asks which such series are algebraic. The record does not give an underlying family of sequences, a desired arithmetic or structural criterion, effective input data, or a required algorithm. The polynomial relation above defines algebraicity but by itself is not an additional characterization of integer sequences. No restriction to a guessed family of special-function, factorial-ratio, or recurrence sequences is inserted. The intended generality and meaning of a satisfactory characterization remain formulation questions.\par The supplied triage checked 2026-08-21 labels the item open but explicitly describes the connection question as incompletely extracted, with no precise theorem to verify or located resolution in its generality. That dated warning is retained, rather than taken as evidence of a particular unformulated conjecture or an independently verified current status.
\subsection{Short English statement}
Describe which integer-coefficient ordinary generating series satisfy a polynomial equation over \(\Q(z)\). The excerpt does not specify a sequence class or the type of structural or arithmetic characterization sought, so the definition of algebraicity alone does not resolve the research request.
\subsection{Sources}
\begin{itemize}\item Local UnsolvedMath registry, permanent ID 30000273, code catalogue entry; \emph{Integral Sequences and Algebraic Generating Functions}, Oberwolfach Reports Open Problems.\item \emph{Explicit Methods in Number Theory (2006)}, Oberwolfach Reports, DOI 10.4171/owr/2005/32, pp. 1799--1866. \url{https://doi.org/10.4171/owr/2005/32}. Proposed year and year tag: 2005; no proposer named.\item Supplied literature label \emph{original Oberwolfach report}: \url{https://doi.org/10.4171/owr/2005/32}. No separate authoritative theorem citation or underlying sequence family is supplied.\item Literature triage checked 2026-08-21, status \texttt{open} with an incomplete-extraction warning; record updated 2026-09-16. The \texttt{corrected\_verified} field concerns notation normalization. Verify the dated assessment independently.\end{itemize}
\end{document}
